
\section{Introduction} \label{sec:introc1}

Financial literacy is at the center of the education debate. As individuals are increasingly exposed to financial decisions at a young age, the sooner they learn to incorporate the costs of their choices into their decision-making, the better. In fact, the Organization for Economic Cooperation and Development (OECD) recommends that financial education start at school, so that people are educated about financial matters as early as possible in their lives \citep[][Good Practice 9]{OECD2005}.\footnote{In this study, we use financial education and financial literacy interchangeably.} In this context, financial education interventions have been implemented in schools around the world, aiming to develop financial skills that could be conducive to better decision-making throughout individuals' life-cycle \citep{Bruhn2016, berry2018impact, Luhrmann2018, gill2019effects, batty2020experiential, de2021effect, frisancho2023school, bover2026impact}.

Most of the aggregate evidence on the impact of financial education, however, comes from programs aimed at adults, and early results were mixed at best: interventions explained a negligible share of the variance in financial behaviors \citep{Fernandes2014} and modestly improved savings behavior but not loan repayment \citep{Miller2015}. More recently, a meta-analysis pooling 76 randomized controlled trials finds that financial education interventions have a significant impact on financial knowledge (about 0.2 of a standard deviation) and on financial behaviors such as saving, borrowing, and budgeting (about 0.1 of a standard deviation) \citep{kaiser2022financial}.\footnote{The meta-analysis covers more than 160,000 individuals; about three quarters of its treatment-effect estimates come from samples of adults, one fifth from youths aged 14 to 25, and fewer than one tenth from children under 14.}

Financial education might change behavior through at least two channels. The first is financial knowledge itself: financially literate individuals are less likely to make costly financial mistakes, such as high-cost borrowing, and are more likely to take advantage of economic opportunities, such as retirement planning and participation in financial markets \citep{Lusardi2014}. In incentivized convex time budget experiments in Germany, a financial education program made high school students' intertemporal choices more time-consistent, though not more patient \citep{Luhrmann2018}, while, absent any intervention, university students with higher pre-existing financial knowledge made seemingly more patient choices, a pattern attributed to knowledge-driven arbitrage rather than to deeper patience \citep{oberrauch2022cognitive}. Nonetheless, whether gains in financial knowledge drive changes in behavior is rarely tested directly.

The second channel is the shaping of attributes such as self-control and the ability to make forward-looking decisions, particularly in programs aimed at children. Such noncognitive attributes shape schooling decisions and labor market prospects \citep{heckman2006effects, kautz2014fostering, acosta2018role}. In particular, more impatient children and adolescents are less likely to invest in human capital and more likely to consume alcohol, develop obesity, and become pregnant during adolescence, while those with lower self-control are more likely to commit crimes later in life \citep{castillo2011today, sutter2013impatience, golsteyn2014adolescent, moffitt2011gradient}. These attributes are also malleable in childhood: a short classroom program, framed as financial literacy, that had young Turkish students imagine their future selves raised their patience in incentivized choices, an effect still detectable almost three years later, and lowered the incidence of behavioral problems at school \citep{alan2018fostering}; a subsequent arm of the same experiment, targeting grit, increased students' perseverance and school performance \citep{alan2019ever}, as did a nationwide grit curriculum in North Macedonia \citep{santos2021can}. In theory, programs aimed at developing such attributes should be prioritized in childhood, given the expected higher returns per dollar invested \citep{cunha2007technology}.

There is increasing evidence that financial education experiments with high school students are effective in increasing financial literacy, with effects ranging from 0.14 to 0.24 of a standard deviation \citep{bover2026impact, frisancho2023school, Bruhn2016}.\footnote{In a quasi-experiment, \cite{gill2019effects} find evidence that twelfth graders taught financial concepts in U.S. economics classes gained about 13 percentage points on a financial literacy test.} The available literature also points to changes in behavior: treated students save and budget more in a large experiment in Brazilian high schools \citep{Bruhn2016}, and these effects can be long-lasting; nine years after the program, treated students used less expensive credit, defaulted less, and were more likely to own a formal microenterprise \citep{bruhn2025long}. In Peru, seven years after school-based lessons, treated individuals shifted their borrowing toward productive small-business loans on slightly better terms, suggesting credit-financed entrepreneurial activity, although formal business creation is unchanged \citep{frisancho2025lasting}. In Spain, treated students make more patient choices in incentivized tasks months after the course ends \citep{bover2026impact}. Evidence from U.S. state financial education mandates shows that students who were required to take personal finance courses in high school have higher credit scores and fewer delinquencies at ages 18 to 21 \citep{urban2020effects}, finance college through federal student aid rather than costlier private loans \citep{stoddard2020effects}, and rely less on payday loans \citep{harvey2019impact}.

There is also growing evidence on the impacts of financial literacy interventions in elementary and middle school grades.
In an experiment in Belgium, \cite{de2021effect} find that a four-hour online course increased the financial knowledge of eighth and ninth graders by 0.46 of a standard deviation,\footnote{Financial knowledge was measured right after the end of the course.} though without a corresponding change in their consumer choices. \cite{batty2020experiential} find an increase of about 0.25 of a standard deviation in the financial knowledge of fourth graders in an experiment with an experiential ``classroom economy'' program across 20 schools in the United States. \cite{berry2018impact}, on the other hand, do not find impacts on financial literacy in an experiment across 135 Ghanaian elementary and junior high schools, with the programs shifting where children saved rather than raising their total savings. Beyond experiments, \cite{bhattacharya2016effectiveness} find a gain of about 12 percentage points in the financial knowledge of eighth graders attending a voluntary summer program in the United States.

In this paper, we use a cluster-randomized controlled trial to investigate the impact of a financial education pilot implemented in 2015 in elementary and middle school grades in two Brazilian municipalities (Manaus and Joinville). We randomly assigned 101 municipal schools to receive the program, and 100 to a control group. Rather than teaching financial facts alone, the intervention aimed to build the attitudes and everyday competencies behind sound financial choices: telling wants from needs, planning and consuming consciously, saving, and weighing the present against the future consequences of one's decisions. The materials anchor these competencies in everyday situations, with students reading real utility bills, comparing cash and installment prices, budgeting school events, and, in the later grades, making decisions inside stories and games that expose common psychological traps, such as ostentation and immediacy. In short, the program sought to help young students become more conscious of the trade-offs embedded in their daily decisions and more forward-looking. The evaluation focused on grades three, five, seven, and nine for reasons discussed later.

We find positive effects on students' financial knowledge of 0.07 of a standard deviation, on average, driven by the middle school grades (0.1 SD). These magnitudes are below the 0.2 of a standard deviation average for financial knowledge in the meta-analysis discussed above, but that average is dominated by adult samples: among children under 14, the pooled estimates carry 95 percent confidence intervals from 0.01 to 0.55 of a standard deviation for financial knowledge and from 0.02 to 0.11 for financial behaviors \citep{kaiser2022financial}. We also find evidence that the program made students' reported attitudes toward consumption and saving more prudent, particularly fifth graders. Futhermore, for the same age cohort, we find suggestive evidence of changes in some contemporary behavioral outcomes, such as talking about money with parents and friends, indicating linkages to short-term behavioral changes. However, we do not detect robust impacts on contemporary learning outcomes or on medium-term human capital. We use a mediation causal analysis to investigate the potential causal link between gains in financial knowledge and changes in behavioral outcomes. Our results do not point to a compelling causal link between formal knowledge and changes in behavior, consistent with \cite{de2021effect}.

Our paper has four main contributions. First, we document the positive impacts on financial proficiency, attitudes, and behavioral outcomes of a large-scale pilot intervention implemented in elementary and middle school grades.  Second, we show that financial education programs can potentially affect different skills across students' life cycles. Third, we test the overlooked hypothesis that increasing financial proficiency is a necessary condition to change individuals' behavior. Finally, we use student-level data from the Brazilian national standardized exams to assess the program's impacts on students' contemporary and medium-term human capital.

Apart from this introduction, this paper is organized as follows. Section \ref{sec:pilot} describes the pilot program and its implementation. Section \ref{sec:data_c} describes the data collection. Section \ref{sec:identification} introduces the empirical strategy, whereas Section \ref{sec:results} discusses the results. We then conclude.

